\chapter{Conclusion}
\label{sec:conclusion}

This report has traced a path from measuring spatial intelligence through improving it to evaluating whether improvements transfer to embodied settings. Along the way, I contributed evaluation infrastructure (nine benchmark integrations, two model integrations, and numerous bug fixes in lmms-eval) that enabled the holistic spatial intelligence evaluation reported in~\cite{cai2025holisticevaluationmultimodalllms} and the model validation in~\cite{cai2025scalingspatialintelligencemultimodal}. I then independently designed and built the EASI embodied evaluation library, a unified toolkit integrating ten embodied benchmarks across five simulators with LLM-powered agents, subprocess isolation for heterogeneous Python environments, and a comprehensive evaluation pipeline with parallel execution, episode filtering, resume support, and trajectory analysis. A preliminary cross-benchmark experiment (EASI-Mini) provided initial data on whether static spatial gains transfer to embodied tasks.

The overarching finding of this research is that spatial intelligence is both more important and more tractable than previously assumed. The gap between models and humans is large but not immutable, and targeted data scaling produces dramatic improvements. The EASI-Mini experiment offers a first look at whether these improvements extend from static benchmarks to embodied competence, and the library provides the infrastructure for the comprehensive studies that will follow.
